\section{Implicaciones metodológicas para fases posteriores}

A partir del estado del arte revisado, una primera implicación para las siguientes fases del proyecto es la conveniencia de comenzar con un \textit{baseline} reproducible y metodológicamente transparente. Una combinación de TF-IDF con un clasificador lineal constituye una referencia adecuada porque permite establecer un punto de comparación claro, barato de entrenar y fácil de auditar antes de escalar hacia modelos más complejos. Este punto de partida debería integrarse dentro de un \textit{pipeline} modular, de modo que etapas como preprocesamiento, representación textual, clasificación y evaluación puedan sustituirse sin rehacer toda la herramienta.

Una segunda implicación es que la validación futura debe priorizar métricas acordes con el riesgo del problema. ROC AUC conviene como métrica central para comparar modelos, pero debe complementarse con \textit{recall}, F1, precisión y un análisis explícito de falsos negativos, dado que omitir casos de riesgo puede ser más costoso que incurrir en algunas falsas alarmas. Asimismo, si el conjunto de datos disponible incluye \texttt{user\_id}, será necesario evitar fuga de información entre entrenamiento y prueba a nivel de usuario, de modo que el desempeño reportado refleje capacidad real de generalización y no reutilización indirecta de patrones del mismo individuo.

Una tercera implicación es que el diseño conceptual de la herramienta deberá incorporar desde el inicio algún mecanismo básico de interpretabilidad. Esto puede materializarse mediante términos relevantes, n-gramas destacados, pesos del clasificador o fragmentos de evidencia que acompañen la predicción. Estas decisiones metodológicas no se formulan aquí como una solución final propuesta, sino como un conjunto de criterios para orientar la Fase 2A; sobre esa base, será posible comparar posteriormente el \textit{baseline} reproducible con variantes más complejas basadas en \textit{transformers} o LLMs.
